\documentclass[11pt]{article}

    \usepackage[breakable]{tcolorbox}
    \usepackage{parskip} % Stop auto-indenting (to mimic markdown behaviour)
    

    % Basic figure setup, for now with no caption control since it's done
    % automatically by Pandoc (which extracts ![](path) syntax from Markdown).
    \usepackage{graphicx}
    % Keep aspect ratio if custom image width or height is specified
    \setkeys{Gin}{keepaspectratio}
    % Maintain compatibility with old templates. Remove in nbconvert 6.0
    \let\Oldincludegraphics\includegraphics
    % Ensure that by default, figures have no caption (until we provide a
    % proper Figure object with a Caption API and a way to capture that
    % in the conversion process - todo).
    \usepackage{caption}
    \DeclareCaptionFormat{nocaption}{}
    \captionsetup{format=nocaption,aboveskip=0pt,belowskip=0pt}

    \usepackage{float}
    \floatplacement{figure}{H} % forces figures to be placed at the correct location
    \usepackage{xcolor} % Allow colors to be defined
    \usepackage{enumerate} % Needed for markdown enumerations to work
    \usepackage{geometry} % Used to adjust the document margins
    \usepackage{amsmath} % Equations
    \usepackage{amssymb} % Equations
    \usepackage{textcomp} % defines textquotesingle
    % Hack from http://tex.stackexchange.com/a/47451/13684:
    \AtBeginDocument{%
        \def\PYZsq{\textquotesingle}% Upright quotes in Pygmentized code
    }
    \usepackage{upquote} % Upright quotes for verbatim code
    \usepackage{eurosym} % defines \euro

    \usepackage{iftex}
    \ifPDFTeX
        \usepackage[T1]{fontenc}
        \IfFileExists{alphabeta.sty}{
              \usepackage{alphabeta}
          }{
              \usepackage[mathletters]{ucs}
              \usepackage[utf8x]{inputenc}
          }
    \else
        \usepackage{fontspec}
        \usepackage{unicode-math}
    \fi

    \usepackage{fancyvrb} % verbatim replacement that allows latex
    \usepackage{grffile} % extends the file name processing of package graphics
                         % to support a larger range
    \makeatletter % fix for old versions of grffile with XeLaTeX
    \@ifpackagelater{grffile}{2019/11/01}
    {
      % Do nothing on new versions
    }
    {
      \def\Gread@@xetex#1{%
        \IfFileExists{"\Gin@base".bb}%
        {\Gread@eps{\Gin@base.bb}}%
        {\Gread@@xetex@aux#1}%
      }
    }
    \makeatother
    \usepackage[Export]{adjustbox} % Used to constrain images to a maximum size
    \adjustboxset{max size={0.9\linewidth}{0.9\paperheight}}

    % The hyperref package gives us a pdf with properly built
    % internal navigation ('pdf bookmarks' for the table of contents,
    % internal cross-reference links, web links for URLs, etc.)
    \usepackage{hyperref}
    % The default LaTeX title has an obnoxious amount of whitespace. By default,
    % titling removes some of it. It also provides customization options.
    \usepackage{titling}
    \usepackage{longtable} % longtable support required by pandoc >1.10
    \usepackage{booktabs}  % table support for pandoc > 1.12.2
    \usepackage{array}     % table support for pandoc >= 2.11.3
    \usepackage{calc}      % table minipage width calculation for pandoc >= 2.11.1
    \usepackage[inline]{enumitem} % IRkernel/repr support (it uses the enumerate* environment)
    \usepackage[normalem]{ulem} % ulem is needed to support strikethroughs (\sout)
                                % normalem makes italics be italics, not underlines
    \usepackage{soul}      % strikethrough (\st) support for pandoc >= 3.0.0
    \usepackage{mathrsfs}
    

    
    % Colors for the hyperref package
    \definecolor{urlcolor}{rgb}{0,.145,.698}
    \definecolor{linkcolor}{rgb}{.71,0.21,0.01}
    \definecolor{citecolor}{rgb}{.12,.54,.11}

    % ANSI colors
    \definecolor{ansi-black}{HTML}{3E424D}
    \definecolor{ansi-black-intense}{HTML}{282C36}
    \definecolor{ansi-red}{HTML}{E75C58}
    \definecolor{ansi-red-intense}{HTML}{B22B31}
    \definecolor{ansi-green}{HTML}{00A250}
    \definecolor{ansi-green-intense}{HTML}{007427}
    \definecolor{ansi-yellow}{HTML}{DDB62B}
    \definecolor{ansi-yellow-intense}{HTML}{B27D12}
    \definecolor{ansi-blue}{HTML}{208FFB}
    \definecolor{ansi-blue-intense}{HTML}{0065CA}
    \definecolor{ansi-magenta}{HTML}{D160C4}
    \definecolor{ansi-magenta-intense}{HTML}{A03196}
    \definecolor{ansi-cyan}{HTML}{60C6C8}
    \definecolor{ansi-cyan-intense}{HTML}{258F8F}
    \definecolor{ansi-white}{HTML}{C5C1B4}
    \definecolor{ansi-white-intense}{HTML}{A1A6B2}
    \definecolor{ansi-default-inverse-fg}{HTML}{FFFFFF}
    \definecolor{ansi-default-inverse-bg}{HTML}{000000}

    % common color for the border for error outputs.
    \definecolor{outerrorbackground}{HTML}{FFDFDF}

    % commands and environments needed by pandoc snippets
    % extracted from the output of `pandoc -s`
    \providecommand{\tightlist}{%
      \setlength{\itemsep}{0pt}\setlength{\parskip}{0pt}}
    \DefineVerbatimEnvironment{Highlighting}{Verbatim}{commandchars=\\\{\}}
    % Add ',fontsize=\small' for more characters per line
    \newenvironment{Shaded}{}{}
    \newcommand{\KeywordTok}[1]{\textcolor[rgb]{0.00,0.44,0.13}{\textbf{{#1}}}}
    \newcommand{\DataTypeTok}[1]{\textcolor[rgb]{0.56,0.13,0.00}{{#1}}}
    \newcommand{\DecValTok}[1]{\textcolor[rgb]{0.25,0.63,0.44}{{#1}}}
    \newcommand{\BaseNTok}[1]{\textcolor[rgb]{0.25,0.63,0.44}{{#1}}}
    \newcommand{\FloatTok}[1]{\textcolor[rgb]{0.25,0.63,0.44}{{#1}}}
    \newcommand{\CharTok}[1]{\textcolor[rgb]{0.25,0.44,0.63}{{#1}}}
    \newcommand{\StringTok}[1]{\textcolor[rgb]{0.25,0.44,0.63}{{#1}}}
    \newcommand{\CommentTok}[1]{\textcolor[rgb]{0.38,0.63,0.69}{\textit{{#1}}}}
    \newcommand{\OtherTok}[1]{\textcolor[rgb]{0.00,0.44,0.13}{{#1}}}
    \newcommand{\AlertTok}[1]{\textcolor[rgb]{1.00,0.00,0.00}{\textbf{{#1}}}}
    \newcommand{\FunctionTok}[1]{\textcolor[rgb]{0.02,0.16,0.49}{{#1}}}
    \newcommand{\RegionMarkerTok}[1]{{#1}}
    \newcommand{\ErrorTok}[1]{\textcolor[rgb]{1.00,0.00,0.00}{\textbf{{#1}}}}
    \newcommand{\NormalTok}[1]{{#1}}

    % Additional commands for more recent versions of Pandoc
    \newcommand{\ConstantTok}[1]{\textcolor[rgb]{0.53,0.00,0.00}{{#1}}}
    \newcommand{\SpecialCharTok}[1]{\textcolor[rgb]{0.25,0.44,0.63}{{#1}}}
    \newcommand{\VerbatimStringTok}[1]{\textcolor[rgb]{0.25,0.44,0.63}{{#1}}}
    \newcommand{\SpecialStringTok}[1]{\textcolor[rgb]{0.73,0.40,0.53}{{#1}}}
    \newcommand{\ImportTok}[1]{{#1}}
    \newcommand{\DocumentationTok}[1]{\textcolor[rgb]{0.73,0.13,0.13}{\textit{{#1}}}}
    \newcommand{\AnnotationTok}[1]{\textcolor[rgb]{0.38,0.63,0.69}{\textbf{\textit{{#1}}}}}
    \newcommand{\CommentVarTok}[1]{\textcolor[rgb]{0.38,0.63,0.69}{\textbf{\textit{{#1}}}}}
    \newcommand{\VariableTok}[1]{\textcolor[rgb]{0.10,0.09,0.49}{{#1}}}
    \newcommand{\ControlFlowTok}[1]{\textcolor[rgb]{0.00,0.44,0.13}{\textbf{{#1}}}}
    \newcommand{\OperatorTok}[1]{\textcolor[rgb]{0.40,0.40,0.40}{{#1}}}
    \newcommand{\BuiltInTok}[1]{{#1}}
    \newcommand{\ExtensionTok}[1]{{#1}}
    \newcommand{\PreprocessorTok}[1]{\textcolor[rgb]{0.74,0.48,0.00}{{#1}}}
    \newcommand{\AttributeTok}[1]{\textcolor[rgb]{0.49,0.56,0.16}{{#1}}}
    \newcommand{\InformationTok}[1]{\textcolor[rgb]{0.38,0.63,0.69}{\textbf{\textit{{#1}}}}}
    \newcommand{\WarningTok}[1]{\textcolor[rgb]{0.38,0.63,0.69}{\textbf{\textit{{#1}}}}}
    \makeatletter
    \newsavebox\pandoc@box
    \newcommand*\pandocbounded[1]{%
      \sbox\pandoc@box{#1}%
      % scaling factors for width and height
      \Gscale@div\@tempa\textheight{\dimexpr\ht\pandoc@box+\dp\pandoc@box\relax}%
      \Gscale@div\@tempb\linewidth{\wd\pandoc@box}%
      % select the smaller of both
      \ifdim\@tempb\p@<\@tempa\p@
        \let\@tempa\@tempb
      \fi
      % scaling accordingly (\@tempa < 1)
      \ifdim\@tempa\p@<\p@
        \scalebox{\@tempa}{\usebox\pandoc@box}%
      % scaling not needed, use as it is
      \else
        \usebox{\pandoc@box}%
      \fi
    }
    \makeatother

    % Define a nice break command that doesn't care if a line doesn't already
    % exist.
    \def\br{\hspace*{\fill} \\* }
    % Math Jax compatibility definitions
    \def\gt{>}
    \def\lt{<}
    \let\Oldtex\TeX
    \let\Oldlatex\LaTeX
    \renewcommand{\TeX}{\textrm{\Oldtex}}
    \renewcommand{\LaTeX}{\textrm{\Oldlatex}}
    % Document parameters
    % Document title
    \title{Kendall Bias}
    
    
    
    
    
    
    
% Pygments definitions
\makeatletter
\def\PY@reset{\let\PY@it=\relax \let\PY@bf=\relax%
    \let\PY@ul=\relax \let\PY@tc=\relax%
    \let\PY@bc=\relax \let\PY@ff=\relax}
\def\PY@tok#1{\csname PY@tok@#1\endcsname}
\def\PY@toks#1+{\ifx\relax#1\empty\else%
    \PY@tok{#1}\expandafter\PY@toks\fi}
\def\PY@do#1{\PY@bc{\PY@tc{\PY@ul{%
    \PY@it{\PY@bf{\PY@ff{#1}}}}}}}
\def\PY#1#2{\PY@reset\PY@toks#1+\relax+\PY@do{#2}}

\@namedef{PY@tok@w}{\def\PY@tc##1{\textcolor[rgb]{0.73,0.73,0.73}{##1}}}
\@namedef{PY@tok@c}{\let\PY@it=\textit\def\PY@tc##1{\textcolor[rgb]{0.24,0.48,0.48}{##1}}}
\@namedef{PY@tok@cp}{\def\PY@tc##1{\textcolor[rgb]{0.61,0.40,0.00}{##1}}}
\@namedef{PY@tok@k}{\let\PY@bf=\textbf\def\PY@tc##1{\textcolor[rgb]{0.00,0.50,0.00}{##1}}}
\@namedef{PY@tok@kp}{\def\PY@tc##1{\textcolor[rgb]{0.00,0.50,0.00}{##1}}}
\@namedef{PY@tok@kt}{\def\PY@tc##1{\textcolor[rgb]{0.69,0.00,0.25}{##1}}}
\@namedef{PY@tok@o}{\def\PY@tc##1{\textcolor[rgb]{0.40,0.40,0.40}{##1}}}
\@namedef{PY@tok@ow}{\let\PY@bf=\textbf\def\PY@tc##1{\textcolor[rgb]{0.67,0.13,1.00}{##1}}}
\@namedef{PY@tok@nb}{\def\PY@tc##1{\textcolor[rgb]{0.00,0.50,0.00}{##1}}}
\@namedef{PY@tok@nf}{\def\PY@tc##1{\textcolor[rgb]{0.00,0.00,1.00}{##1}}}
\@namedef{PY@tok@nc}{\let\PY@bf=\textbf\def\PY@tc##1{\textcolor[rgb]{0.00,0.00,1.00}{##1}}}
\@namedef{PY@tok@nn}{\let\PY@bf=\textbf\def\PY@tc##1{\textcolor[rgb]{0.00,0.00,1.00}{##1}}}
\@namedef{PY@tok@ne}{\let\PY@bf=\textbf\def\PY@tc##1{\textcolor[rgb]{0.80,0.25,0.22}{##1}}}
\@namedef{PY@tok@nv}{\def\PY@tc##1{\textcolor[rgb]{0.10,0.09,0.49}{##1}}}
\@namedef{PY@tok@no}{\def\PY@tc##1{\textcolor[rgb]{0.53,0.00,0.00}{##1}}}
\@namedef{PY@tok@nl}{\def\PY@tc##1{\textcolor[rgb]{0.46,0.46,0.00}{##1}}}
\@namedef{PY@tok@ni}{\let\PY@bf=\textbf\def\PY@tc##1{\textcolor[rgb]{0.44,0.44,0.44}{##1}}}
\@namedef{PY@tok@na}{\def\PY@tc##1{\textcolor[rgb]{0.41,0.47,0.13}{##1}}}
\@namedef{PY@tok@nt}{\let\PY@bf=\textbf\def\PY@tc##1{\textcolor[rgb]{0.00,0.50,0.00}{##1}}}
\@namedef{PY@tok@nd}{\def\PY@tc##1{\textcolor[rgb]{0.67,0.13,1.00}{##1}}}
\@namedef{PY@tok@s}{\def\PY@tc##1{\textcolor[rgb]{0.73,0.13,0.13}{##1}}}
\@namedef{PY@tok@sd}{\let\PY@it=\textit\def\PY@tc##1{\textcolor[rgb]{0.73,0.13,0.13}{##1}}}
\@namedef{PY@tok@si}{\let\PY@bf=\textbf\def\PY@tc##1{\textcolor[rgb]{0.64,0.35,0.47}{##1}}}
\@namedef{PY@tok@se}{\let\PY@bf=\textbf\def\PY@tc##1{\textcolor[rgb]{0.67,0.36,0.12}{##1}}}
\@namedef{PY@tok@sr}{\def\PY@tc##1{\textcolor[rgb]{0.64,0.35,0.47}{##1}}}
\@namedef{PY@tok@ss}{\def\PY@tc##1{\textcolor[rgb]{0.10,0.09,0.49}{##1}}}
\@namedef{PY@tok@sx}{\def\PY@tc##1{\textcolor[rgb]{0.00,0.50,0.00}{##1}}}
\@namedef{PY@tok@m}{\def\PY@tc##1{\textcolor[rgb]{0.40,0.40,0.40}{##1}}}
\@namedef{PY@tok@gh}{\let\PY@bf=\textbf\def\PY@tc##1{\textcolor[rgb]{0.00,0.00,0.50}{##1}}}
\@namedef{PY@tok@gu}{\let\PY@bf=\textbf\def\PY@tc##1{\textcolor[rgb]{0.50,0.00,0.50}{##1}}}
\@namedef{PY@tok@gd}{\def\PY@tc##1{\textcolor[rgb]{0.63,0.00,0.00}{##1}}}
\@namedef{PY@tok@gi}{\def\PY@tc##1{\textcolor[rgb]{0.00,0.52,0.00}{##1}}}
\@namedef{PY@tok@gr}{\def\PY@tc##1{\textcolor[rgb]{0.89,0.00,0.00}{##1}}}
\@namedef{PY@tok@ge}{\let\PY@it=\textit}
\@namedef{PY@tok@gs}{\let\PY@bf=\textbf}
\@namedef{PY@tok@ges}{\let\PY@bf=\textbf\let\PY@it=\textit}
\@namedef{PY@tok@gp}{\let\PY@bf=\textbf\def\PY@tc##1{\textcolor[rgb]{0.00,0.00,0.50}{##1}}}
\@namedef{PY@tok@go}{\def\PY@tc##1{\textcolor[rgb]{0.44,0.44,0.44}{##1}}}
\@namedef{PY@tok@gt}{\def\PY@tc##1{\textcolor[rgb]{0.00,0.27,0.87}{##1}}}
\@namedef{PY@tok@err}{\def\PY@bc##1{{\setlength{\fboxsep}{\string -\fboxrule}\fcolorbox[rgb]{1.00,0.00,0.00}{1,1,1}{\strut ##1}}}}
\@namedef{PY@tok@kc}{\let\PY@bf=\textbf\def\PY@tc##1{\textcolor[rgb]{0.00,0.50,0.00}{##1}}}
\@namedef{PY@tok@kd}{\let\PY@bf=\textbf\def\PY@tc##1{\textcolor[rgb]{0.00,0.50,0.00}{##1}}}
\@namedef{PY@tok@kn}{\let\PY@bf=\textbf\def\PY@tc##1{\textcolor[rgb]{0.00,0.50,0.00}{##1}}}
\@namedef{PY@tok@kr}{\let\PY@bf=\textbf\def\PY@tc##1{\textcolor[rgb]{0.00,0.50,0.00}{##1}}}
\@namedef{PY@tok@bp}{\def\PY@tc##1{\textcolor[rgb]{0.00,0.50,0.00}{##1}}}
\@namedef{PY@tok@fm}{\def\PY@tc##1{\textcolor[rgb]{0.00,0.00,1.00}{##1}}}
\@namedef{PY@tok@vc}{\def\PY@tc##1{\textcolor[rgb]{0.10,0.09,0.49}{##1}}}
\@namedef{PY@tok@vg}{\def\PY@tc##1{\textcolor[rgb]{0.10,0.09,0.49}{##1}}}
\@namedef{PY@tok@vi}{\def\PY@tc##1{\textcolor[rgb]{0.10,0.09,0.49}{##1}}}
\@namedef{PY@tok@vm}{\def\PY@tc##1{\textcolor[rgb]{0.10,0.09,0.49}{##1}}}
\@namedef{PY@tok@sa}{\def\PY@tc##1{\textcolor[rgb]{0.73,0.13,0.13}{##1}}}
\@namedef{PY@tok@sb}{\def\PY@tc##1{\textcolor[rgb]{0.73,0.13,0.13}{##1}}}
\@namedef{PY@tok@sc}{\def\PY@tc##1{\textcolor[rgb]{0.73,0.13,0.13}{##1}}}
\@namedef{PY@tok@dl}{\def\PY@tc##1{\textcolor[rgb]{0.73,0.13,0.13}{##1}}}
\@namedef{PY@tok@s2}{\def\PY@tc##1{\textcolor[rgb]{0.73,0.13,0.13}{##1}}}
\@namedef{PY@tok@sh}{\def\PY@tc##1{\textcolor[rgb]{0.73,0.13,0.13}{##1}}}
\@namedef{PY@tok@s1}{\def\PY@tc##1{\textcolor[rgb]{0.73,0.13,0.13}{##1}}}
\@namedef{PY@tok@mb}{\def\PY@tc##1{\textcolor[rgb]{0.40,0.40,0.40}{##1}}}
\@namedef{PY@tok@mf}{\def\PY@tc##1{\textcolor[rgb]{0.40,0.40,0.40}{##1}}}
\@namedef{PY@tok@mh}{\def\PY@tc##1{\textcolor[rgb]{0.40,0.40,0.40}{##1}}}
\@namedef{PY@tok@mi}{\def\PY@tc##1{\textcolor[rgb]{0.40,0.40,0.40}{##1}}}
\@namedef{PY@tok@il}{\def\PY@tc##1{\textcolor[rgb]{0.40,0.40,0.40}{##1}}}
\@namedef{PY@tok@mo}{\def\PY@tc##1{\textcolor[rgb]{0.40,0.40,0.40}{##1}}}
\@namedef{PY@tok@ch}{\let\PY@it=\textit\def\PY@tc##1{\textcolor[rgb]{0.24,0.48,0.48}{##1}}}
\@namedef{PY@tok@cm}{\let\PY@it=\textit\def\PY@tc##1{\textcolor[rgb]{0.24,0.48,0.48}{##1}}}
\@namedef{PY@tok@cpf}{\let\PY@it=\textit\def\PY@tc##1{\textcolor[rgb]{0.24,0.48,0.48}{##1}}}
\@namedef{PY@tok@c1}{\let\PY@it=\textit\def\PY@tc##1{\textcolor[rgb]{0.24,0.48,0.48}{##1}}}
\@namedef{PY@tok@cs}{\let\PY@it=\textit\def\PY@tc##1{\textcolor[rgb]{0.24,0.48,0.48}{##1}}}

\def\PYZbs{\char`\\}
\def\PYZus{\char`\_}
\def\PYZob{\char`\{}
\def\PYZcb{\char`\}}
\def\PYZca{\char`\^}
\def\PYZam{\char`\&}
\def\PYZlt{\char`\<}
\def\PYZgt{\char`\>}
\def\PYZsh{\char`\#}
\def\PYZpc{\char`\%}
\def\PYZdl{\char`\$}
\def\PYZhy{\char`\-}
\def\PYZsq{\char`\'}
\def\PYZdq{\char`\"}
\def\PYZti{\char`\~}
% for compatibility with earlier versions
\def\PYZat{@}
\def\PYZlb{[}
\def\PYZrb{]}
\makeatother


    % For linebreaks inside Verbatim environment from package fancyvrb.
    \makeatletter
        \newbox\Wrappedcontinuationbox
        \newbox\Wrappedvisiblespacebox
        \newcommand*\Wrappedvisiblespace {\textcolor{red}{\textvisiblespace}}
        \newcommand*\Wrappedcontinuationsymbol {\textcolor{red}{\llap{\tiny$\m@th\hookrightarrow$}}}
        \newcommand*\Wrappedcontinuationindent {3ex }
        \newcommand*\Wrappedafterbreak {\kern\Wrappedcontinuationindent\copy\Wrappedcontinuationbox}
        % Take advantage of the already applied Pygments mark-up to insert
        % potential linebreaks for TeX processing.
        %        {, <, #, %, $, ' and ": go to next line.
        %        _, }, ^, &, >, - and ~: stay at end of broken line.
        % Use of \textquotesingle for straight quote.
        \newcommand*\Wrappedbreaksatspecials {%
            \def\PYGZus{\discretionary{\char`\_}{\Wrappedafterbreak}{\char`\_}}%
            \def\PYGZob{\discretionary{}{\Wrappedafterbreak\char`\{}{\char`\{}}%
            \def\PYGZcb{\discretionary{\char`\}}{\Wrappedafterbreak}{\char`\}}}%
            \def\PYGZca{\discretionary{\char`\^}{\Wrappedafterbreak}{\char`\^}}%
            \def\PYGZam{\discretionary{\char`\&}{\Wrappedafterbreak}{\char`\&}}%
            \def\PYGZlt{\discretionary{}{\Wrappedafterbreak\char`\<}{\char`\<}}%
            \def\PYGZgt{\discretionary{\char`\>}{\Wrappedafterbreak}{\char`\>}}%
            \def\PYGZsh{\discretionary{}{\Wrappedafterbreak\char`\#}{\char`\#}}%
            \def\PYGZpc{\discretionary{}{\Wrappedafterbreak\char`\%}{\char`\%}}%
            \def\PYGZdl{\discretionary{}{\Wrappedafterbreak\char`\$}{\char`\$}}%
            \def\PYGZhy{\discretionary{\char`\-}{\Wrappedafterbreak}{\char`\-}}%
            \def\PYGZsq{\discretionary{}{\Wrappedafterbreak\textquotesingle}{\textquotesingle}}%
            \def\PYGZdq{\discretionary{}{\Wrappedafterbreak\char`\"}{\char`\"}}%
            \def\PYGZti{\discretionary{\char`\~}{\Wrappedafterbreak}{\char`\~}}%
        }
        % Some characters . , ; ? ! / are not pygmentized.
        % This macro makes them "active" and they will insert potential linebreaks
        \newcommand*\Wrappedbreaksatpunct {%
            \lccode`\~`\.\lowercase{\def~}{\discretionary{\hbox{\char`\.}}{\Wrappedafterbreak}{\hbox{\char`\.}}}%
            \lccode`\~`\,\lowercase{\def~}{\discretionary{\hbox{\char`\,}}{\Wrappedafterbreak}{\hbox{\char`\,}}}%
            \lccode`\~`\;\lowercase{\def~}{\discretionary{\hbox{\char`\;}}{\Wrappedafterbreak}{\hbox{\char`\;}}}%
            \lccode`\~`\:\lowercase{\def~}{\discretionary{\hbox{\char`\:}}{\Wrappedafterbreak}{\hbox{\char`\:}}}%
            \lccode`\~`\?\lowercase{\def~}{\discretionary{\hbox{\char`\?}}{\Wrappedafterbreak}{\hbox{\char`\?}}}%
            \lccode`\~`\!\lowercase{\def~}{\discretionary{\hbox{\char`\!}}{\Wrappedafterbreak}{\hbox{\char`\!}}}%
            \lccode`\~`\/\lowercase{\def~}{\discretionary{\hbox{\char`\/}}{\Wrappedafterbreak}{\hbox{\char`\/}}}%
            \catcode`\.\active
            \catcode`\,\active
            \catcode`\;\active
            \catcode`\:\active
            \catcode`\?\active
            \catcode`\!\active
            \catcode`\/\active
            \lccode`\~`\~
        }
    \makeatother

    \let\OriginalVerbatim=\Verbatim
    \makeatletter
    \renewcommand{\Verbatim}[1][1]{%
        %\parskip\z@skip
        \sbox\Wrappedcontinuationbox {\Wrappedcontinuationsymbol}%
        \sbox\Wrappedvisiblespacebox {\FV@SetupFont\Wrappedvisiblespace}%
        \def\FancyVerbFormatLine ##1{\hsize\linewidth
            \vtop{\raggedright\hyphenpenalty\z@\exhyphenpenalty\z@
                \doublehyphendemerits\z@\finalhyphendemerits\z@
                \strut ##1\strut}%
        }%
        % If the linebreak is at a space, the latter will be displayed as visible
        % space at end of first line, and a continuation symbol starts next line.
        % Stretch/shrink are however usually zero for typewriter font.
        \def\FV@Space {%
            \nobreak\hskip\z@ plus\fontdimen3\font minus\fontdimen4\font
            \discretionary{\copy\Wrappedvisiblespacebox}{\Wrappedafterbreak}
            {\kern\fontdimen2\font}%
        }%

        % Allow breaks at special characters using \PYG... macros.
        \Wrappedbreaksatspecials
        % Breaks at punctuation characters . , ; ? ! and / need catcode=\active
        \OriginalVerbatim[#1,codes*=\Wrappedbreaksatpunct]%
    }
    \makeatother

    % Exact colors from NB
    \definecolor{incolor}{HTML}{303F9F}
    \definecolor{outcolor}{HTML}{D84315}
    \definecolor{cellborder}{HTML}{CFCFCF}
    \definecolor{cellbackground}{HTML}{F7F7F7}

    % prompt
    \makeatletter
    \newcommand{\boxspacing}{\kern\kvtcb@left@rule\kern\kvtcb@boxsep}
    \makeatother
    \newcommand{\prompt}[4]{
        {\ttfamily\llap{{\color{#2}[#3]:\hspace{3pt}#4}}\vspace{-\baselineskip}}
    }
    

    
    % Prevent overflowing lines due to hard-to-break entities
    \sloppy
    % Setup hyperref package
    \hypersetup{
      breaklinks=true,  % so long urls are correctly broken across lines
      colorlinks=true,
      urlcolor=urlcolor,
      linkcolor=linkcolor,
      citecolor=citecolor,
      }
    % Slightly bigger margins than the latex defaults
    
    \geometry{verbose,tmargin=1in,bmargin=1in,lmargin=1in,rmargin=1in}
    
    

\begin{document}
    
    \maketitle
    
    

    
    \section{Problem setup}\label{problem-setup}

We consider the AR(1) data generating process \[
y_t = \rho y_{t-1}+e_t
\] where \(\{e_t\}\) is a white noise process with mean zero and
variance \(\sigma^2\). The goal is to estimate the persistence parameter
\(\rho\) using ordinary least squares.

We write the same model in regression form \[
y_t = \beta_1 y_{t-1}+u_t
\] where \(\beta_1\) is the parameter to be estimated and \(u_t\) is the
regression error term. Under the true data generating process, we have
\(\beta_1 = \rho\) and \(u_t = e_t\).

The OLS estimator for \(\beta_1\) is \[
\begin{aligned}
\hat\beta_{1} &= \frac{\sum_{t=2}^{T} y_t y_{t-1}}{\sum_{t=2}^{T} y_{t-1}^2}
\end{aligned}
\]

Substituting the true model \(y_t = \rho y_{t-1} + e_t\) into the
numerator gives \[
\begin{aligned}
\hat\beta_{1}
&= \frac{\sum_{t=2}^{T} (\rho y_{t-1}+e_t) y_{t-1}}{\sum_{t=2}^{T} y_{t-1}^2} \\
&= \frac{\rho \sum_{t=2}^{T} y_{t-1}^2 + \sum_{t=2}^{T} e_t y_{t-1}}{\sum_{t=2}^{T} y_{t-1}^2} \\
&= \rho+\frac{\sum_{t=2}^{T} e_t y_{t-1}}{\sum_{t=2}^{T} y_{t-1}^2}
\end{aligned}
\]

This decomposition is the key starting point for the bias calculation.
It shows that the estimator differs from the true parameter \(\rho\) by
a ratio of two random quantities: - the numerator \(\sum e_t y_{t-1}\),
which captures the correlation between the error and the regressor - the
denominator \(\sum y_{t-1}^2\), which rescales the estimator

Even though \(e_t\) is uncorrelated with \(y_{t-1}\) in expectation at a
fixed time, the ratio structure introduces bias in finite samples. The
rest of the derivation focuses on understanding the expectation of this
ratio.

\subsection{Background computations}\label{background-computations}

To analyze the terms above, we first express \(y_t\) in terms of initial
conditions and the sequence of shocks.

Starting from the AR(1) recursion: \[
\begin{aligned}
y_t&=\rho y_{t-1}+e_t\\
&=\rho^2 y_{t-2}+\rho e_{t-1}+e_t\\
&=\rho^t y_{0}+\sum_{s=0}^{t-1}\rho^s e_{t-s}
\end{aligned}
\]

This representation writes \(y_t\) as a weighted sum of past shocks plus
the initial condition. It makes clear that \(y_t\) depends on the entire
history of innovations, with weights that decay geometrically at rate
\(\rho\).

We next compute the second moment of \(y_t\): \[
\text{E}\left[y_{t}^2\right]
\]

Using the moving average representation and the independence of the
shocks, cross terms vanish, leaving \[
\text{E}\left[y_{t}^2\right]
=\rho^{2t} \text{E}\left[y_0^2\right]+\sum_{s=0}^{t-1}\rho^{2s}\text{E}\left[e_{t-s}^2\right]
\]

Assuming stationarity, we set
\(\text{E}[y_0^2] = \frac{\sigma^2}{1-\rho^2}\) and
\(\text{E}[e_t^2] = \sigma^2\), giving \[
=\rho^{2t} \frac{\sigma^2}{1-\rho^2}+\sigma^2\sum_{s=0}^{t-1}\rho^{2s}
\]

The geometric sum simplifies to \[
\sum_{s=0}^{t-1}\rho^{2s} = \frac{1 - \rho^{2t}}{1 - \rho^2}
\]

Substituting this in, the dependence on \(t\) cancels out, yielding \[
\text{E}\left[y_t^2\right] = \frac{\sigma^2}{1-\rho^2}
\]

This shows that the variance of \(y_t\) is constant over time under
stationarity. This result will be used repeatedly when evaluating
expectations involving the denominator \(\sum y_{t-1}^2\).

    \section{Taylor expansion framework}\label{taylor-expansion-framework}

From the previous decomposition, we write the estimation error as \[
\hat\beta_{1}-\rho = \frac{(T-1)^{-1}\sum_{t=2}^{T} e_t y_{t-1}}{(T-1)^{-1}\sum_{t=2}^{T} y_{t-1}^2} = \frac{A}{B}
\]

where we define the sample averages -
\(A = (T-1)^{-1}\sum_{t=2}^{T} e_t y_{t-1}\) -
\(B = (T-1)^{-1}\sum_{t=2}^{T} y_{t-1}^2\)

The goal is to compute
\(\text{E}[\hat\beta_1 - \rho] = \text{E}[A/B]\).\\
The difficulty is that this is the expectation of a ratio of random
variables, which does not simplify to the ratio of expectations.

To handle this, we approximate the ratio using a Taylor expansion around
the means of \(A\) and \(B\). Let \[
a = \text{E}[A], \quad b = \text{E}[B]
\]

We expand the function \(f(A,B) = A/B\) around the point \((a,b)\) up to
second order: \[
\frac{A}{B}=\frac{a}{b}+\frac{1}{B}(A-a)-\frac{a}{b^2}(B-b)-\frac{1}{b^2}(A-a)(B-b)+\frac{a}{b^3}(B-b)^2+R_2
\]

This expansion expresses the ratio as: - a deterministic leading term
\(\frac{a}{b}\) - linear corrections in \((A-a)\) and \((B-b)\) -
second-order interaction terms - a remainder \(R_2\) capturing
higher-order terms

Taking expectations term by term, we obtain \[
\begin{aligned}
\text{E}\left[\frac{A}{B}\right]
&=\frac{a}{b}
+\frac{1}{B}\text{E}\left[A-a\right]
-\frac{a}{b^2}\text{E}\left[B-b\right] \\
&\quad -\frac{1}{b^2}\text{E}\left[(A-a)(B-b)\right]
+\frac{a}{b^3}\text{E}\left[(B-b)^2\right]
+\text{E}\left[R_2\right]
\end{aligned}
\]

Now we simplify each term: - \(\text{E}[A-a] = 0\) and
\(\text{E}[B-b] = 0\) by definition -
\(\text{E}[(A-a)(B-b)] = \text{Cov}(A,B)\) -
\(\text{E}[(B-b)^2] = \text{Var}(B)\)

Substituting these in gives \[
\text{E}\left[\frac{A}{B}\right]
=\frac{a}{b}
-\frac{1}{b^2}\text{Cov}\left[A,B\right]
+\frac{a}{b^3}\text{Var}\left[B\right]
+\text{E}\left[R_2\right]
\]

Finally, rewriting everything in terms of expectations, \[
\text{E}\left[\frac{A}{B}\right]
=\frac{\text{E}[A]}{\text{E}[B]}
-\frac{1}{\text{E}[B]^2}\text{Cov}\left[A,B\right]
+\frac{\text{E}[A]}{\text{E}[B]^3}\text{Var}\left[B\right]
+\text{E}\left[R_2\right]
\]

This expression is the central tool for deriving the Kendall bias. It
shows that the bias arises from two key sources: - the covariance
between numerator and denominator, \(\text{Cov}(A,B)\) - the variability
of the denominator, \(\text{Var}(B)\)

Even if \(\text{E}[A] = 0\), the estimator can still be biased due to
these second-order terms. The remainder term \(\text{E}[R_2]\) is of
smaller order and will typically vanish asymptotically, allowing us to
focus on the leading bias contributions.

    \section{Warm-up: No Constant case}\label{warm-up-no-constant-case}

In this section, we compute the basic moments of \(A\) and \(B\) in the
simpler case where there is no intercept in the model. This serves as a
warm-up before tackling the full bias expression.

Recall the definitions - \(A = (T-1)^{-1}\sum_{t=2}^{T} e_t y_{t-1}\) -
\(B = (T-1)^{-1}\sum_{t=2}^{T} y_{t-1}^2\)

Our goal is to compute \(\text{E}[A]\) and \(\text{E}[B]\), which will
be used in the Taylor expansion.

\subsection{Definitions}\label{definitions}

We begin with the expectation of \(A\): \[
\begin{aligned}
\text{E}\left[A\right]
&=\frac{1}{T-1}\sum_{t=2}^{T} \text{E}\left[e_t y_{t-1}\right]
\end{aligned}
\]

Because \(e_t\) is independent of past values (in particular
\(y_{t-1}\)), we can factor the expectation: \[
\begin{aligned}
&=\frac{1}{T-1}\sum_{t=2}^{T} \text{E}\left[e_t\right]\text{E}\left[y_{t-1}\right]
\end{aligned}
\]

Since \(\text{E}[e_t]=0\) and (in the no-constant case)
\(\text{E}[y_{t-1}]=0\), each term in the sum is zero, giving \[
\text{E}\left[A\right]=0
\]

This is an important baseline result: on average, the numerator of the
estimator is zero. As we will see, the bias does not come from
\(\text{E}[A]\), but from higher-order interactions between \(A\) and
\(B\).

\begin{center}\rule{0.5\linewidth}{0.5pt}\end{center}

Next, we compute the expectation of \(B\): \[
\begin{aligned}
\text{E}\left[B\right]
&=\frac{1}{T-1}\sum_{t=2}^{T} \text{E}\left[y_{t-1}^2\right]
\end{aligned}
\]

From the background computations, we know that under stationarity \[
\text{E}\left[y_{t}^2\right] = \frac{\sigma^2}{1-\rho^{2}}
\]

This does not depend on \(t\), so every term in the sum is identical: \[
\begin{aligned}
\text{E}\left[B\right]
&=\frac{1}{T-1}\sum_{t=2}^{T} \frac{\sigma^2}{1-\rho^{2}}\\
&=\frac{\sigma^2}{1-\rho^{2}}
\end{aligned}
\]

Thus, the denominator converges to the unconditional variance of the
AR(1) process.

\begin{center}\rule{0.5\linewidth}{0.5pt}\end{center}

For completeness, we restate the calculation of \(\text{E}[y_t^2]\),
which underlies the result above: \[
\begin{aligned}
\text{E}\left[y_{t}^2\right]
&=\rho^{2t} \sigma_0+\sigma^2\frac{1-\rho^{2t}}{1-\rho^{2}}
\end{aligned}
\]

Substituting the stationary initial variance
\(\sigma_0 = \frac{\sigma^2}{1-\rho^2}\) gives \[
\begin{aligned}
&=\rho^{2t} \frac{\sigma^2}{1-\rho^{2}}+\sigma^2\frac{1-\rho^{2t}}{1-\rho^{2}}
\end{aligned}
\]

Combining terms, the dependence on \(t\) cancels: \[
\text{E}\left[y_{t}^2\right]
=\frac{\sigma^2}{1-\rho^{2}}
\]

This confirms that the process has constant variance over time, which
justifies replacing each term in \(\text{E}[B]\) with the same value.

\begin{center}\rule{0.5\linewidth}{0.5pt}\end{center}

At this point, we have established: - \(\text{E}[A] = 0\) -
\(\text{E}[B] = \frac{\sigma^2}{1-\rho^2}\)

These results will simplify the Taylor expansion considerably, since
many terms involving \(\text{E}[A]\) will vanish.

    \subsection{\texorpdfstring{Compute
\(\text{Cov}[A,B]\)}{Compute \textbackslash text\{Cov\}{[}A,B{]}}}\label{compute-textcovab}

We now compute the key second-order term driving the bias: \[
\text{Cov}[A,B]
\]

Recall that - \(A = (T-1)^{-1}\sum_{t=2}^T e_t y_{t-1}\) -
\(B = (T-1)^{-1}\sum_{t=2}^T y_{t-1}^2\)

Since we have already shown that \(\text{E}[A] = 0\) and
\(\text{E}[B] = (1-\rho^2)^{-1}\sigma^2\), the covariance simplifies to
\[
\begin{aligned}
\text{Cov}[A,B]
&=\text{E}[(A-\text{E}[A])(B-\text{E}[B])]\\
&=\text{E}\left[A\left(B-(1-\rho^2)^{-1}\sigma^2\right)\right]
\end{aligned}
\]

Substituting the definitions of \(A\) and \(B\), we obtain \[
\begin{aligned}
\text{Cov}[A,B]
&=\text{E}\left[(T-1)^{-2}\sum_{t=2}^T\sum_{s=2}^T e_t y_{t-1}y_{s-1}^2\right]
\end{aligned}
\]

This is a double sum over all pairs \((t,s)\). The next step is to
determine which terms have nonzero expectation.

\begin{center}\rule{0.5\linewidth}{0.5pt}\end{center}

\subsubsection{Step 1: Use independence to restrict the
summation}\label{step-1-use-independence-to-restrict-the-summation}

Because \(e_t\) is independent of all past shocks, terms where
\(s \le t\) involve \(y_{s-1}\) that depends only on shocks up to time
\(s-1 \le t-1\). These are independent of \(e_t\), so their expectation
is zero.

Thus, only terms with \(s > t\) contribute: \[
\begin{aligned}
\text{Cov}[A,B]
&=\text{E}\left[(T-1)^{-2}\sum_{t=2}^T\sum_{s=t+1}^T e_t y_{t-1}y_{s-1}^2\right]
\end{aligned}
\]

This is the key structural simplification: only ``forward-looking''
interactions matter.

\begin{center}\rule{0.5\linewidth}{0.5pt}\end{center}

\subsubsection{Step 2: Expand using the AR(1)
representation}\label{step-2-expand-using-the-ar1-representation}

To evaluate the expectation, we express \(y_{t-1}\) and \(y_{s-1}\) in
terms of shocks: \[
y_{t-1}=\rho^{t-1}y_{0}+\sum_{r=0}^{t-2}\rho^r e_{t-1-r},
\qquad
y_{s-1}=\rho^{s-1}y_{0}+\sum_{q=0}^{s-2}\rho^q e_{s-1-q}
\]

Substituting these into \(e_t y_{t-1} y_{s-1}^2\) produces many terms
involving products of \(y_0\) and various shocks.

After expanding, we group terms by the number of shocks involved: \[
\begin{aligned}
\text{Cov}[A,B]
&=\text{E}\Bigg[
(T-1)^{-2}\sum_{t=2}^T\sum_{s=t+1}^T
\rho^{t-1}\rho^{2(s-1)} y_0^3\, e_t \\
&\quad+
(T-1)^{-2}\sum_{t=2}^T\sum_{s=t+1}^T\sum_{r=0}^{t-2}
\rho^{r}\rho^{2(s-1)} y_0^2\, e_t e_{t-1-r} \\
&\quad+
(T-1)^{-2}\sum_{t=2}^T\sum_{s=t+1}^T\sum_{p=0}^{s-2}
2\rho^{t-1}\rho^{s-1}\rho^{p} y_0^2\, e_t e_{s-1-p} \\
&\quad+
(T-1)^{-2}\sum_{t=2}^T\sum_{s=t+1}^T\sum_{r,q,p}
\rho^{r}\rho^{q+p}\, e_t e_{t-1-r} e_{s-1-q} e_{s-1-p}
\Bigg]
\end{aligned}
\]

\begin{center}\rule{0.5\linewidth}{0.5pt}\end{center}

\subsubsection{Step 3: Eliminate terms with zero
expectation}\label{step-3-eliminate-terms-with-zero-expectation}

Because all shocks have mean zero, any term involving an odd number of
shocks has expectation zero.

This immediately removes: - all terms involving \(y_0^3 e_t\) - all
terms involving three shocks

The only surviving contributions are those involving \textbf{even
numbers of shocks}, giving: \[
\begin{aligned}
\text{Cov}[A,B]
&=
(T-1)^{-2}\sum_{t=2}^T\sum_{s=t+1}^T\sum_{p=0}^{s-2}
2\rho^{t-1}\rho^{s-1}\rho^{p}\,\text{E}[y_0^2 e_t e_{s-1-p}]\\
&\quad+
(T-1)^{-2}\sum_{t=2}^T\sum_{s=t+1}^T\sum_{r=0}^{t-2}\sum_{q,p}
\rho^{r}\rho^{q+p}\,\text{E}[e_t e_{t-1-r} e_{s-1-q} e_{s-1-p}]
\end{aligned}
\]

\begin{center}\rule{0.5\linewidth}{0.5pt}\end{center}

\subsubsection{Step 4: Use independence and index
matching}\label{step-4-use-independence-and-index-matching}

Now we evaluate expectations by matching indices: -
\(\text{E}[e_t e_{s-1-p}] \neq 0\) only if \(t = s-1-p\) - fourth
moments factor into products of second moments under independence

Using these rules, the expression simplifies to \[
\begin{aligned}
\text{Cov}[A,B]
&=
(T-1)^{-2}\sum_{t=2}^T\sum_{s=t+1}^T
2\rho^{2s-3}(1-\rho^2)^{-1}\sigma^4\\
&\quad+
(T-1)^{-2}\sum_{t=2}^T\sum_{s=t+1}^T\sum_{r=0}^{t-2}
2\rho^{2r+2s-2t-1}\sigma^4
\end{aligned}
\]

The first term comes from interactions involving \(y_0^2\), while the
second comes from purely shock-based terms.

\begin{center}\rule{0.5\linewidth}{0.5pt}\end{center}

\subsubsection{Step 5: Evaluate the geometric
sum}\label{step-5-evaluate-the-geometric-sum}

The inner sum over \(r\) is geometric: \[
\begin{aligned}
\sum_{r=0}^{t-2}\rho^{2r+2s-2t-1}
=
\rho^{2s-2t-1}\frac{1-\rho^{2(t-1)}}{1-\rho^2}
\end{aligned}
\]

Substituting this back gives \[
\begin{aligned}
\text{Cov}[A,B]
&=
2(T-1)^{-2}(1-\rho^2)^{-1}\sigma^4
\left[
\sum_{t=2}^T\sum_{s=t+1}^T\rho^{2s-3}
+
\sum_{t=2}^T\sum_{s=t+1}^T\rho^{2s-2t-1}(1-\rho^{2(t-1)})
\right]
\end{aligned}
\]

\begin{center}\rule{0.5\linewidth}{0.5pt}\end{center}

\subsubsection{Step 6: Final
simplification}\label{step-6-final-simplification}

After evaluating the remaining sums (again using geometric series), we
obtain the closed form: \[
\begin{aligned}
\text{Cov}\left[A,B\right]
&=
\frac{2\sigma^4\left((T-1)(\rho^4-\rho^2)-\rho^{2T}\right)}
{\rho (T-1)^2(1-\rho^2)^3}
\end{aligned}
\]

For large \(T\), the term involving \(\rho^{2T}\) vanishes, giving the
leading-order approximation: \[
\text{Cov}\left[A,B\right]
=\frac{2\sigma^4\left(\rho^4-\rho^2\right)}
{\rho (T-1)(1-\rho^2)^3}+O(T^{-2})
\]

\begin{center}\rule{0.5\linewidth}{0.5pt}\end{center}

\subsubsection{Interpretation}\label{interpretation}

This result shows that \(\text{Cov}[A,B]\) is of order \(O(1/T)\).\\
This is exactly the scale at which the Kendall bias appears.

Importantly: - the covariance is \textbf{not zero}, even though
\(\text{E}[A]=0\) - this nonzero covariance is the dominant source of
bias in the estimator

This term will feed directly into the Taylor expansion, producing the
familiar finite-sample bias of the AR(1) OLS estimator.

    \subsection{Complete setup}\label{complete-setup}

We now combine all previous results to obtain the bias of the estimator.

Recall the Taylor expansion for the expectation of the ratio: \[
\text{E}\left[\frac{A}{B}\right]
=\frac{\text{E}[A]}{\text{E}[B]}
-\frac{1}{\text{E}[B]^2}\text{Cov}\left[A,B\right]
+\frac{\text{E}[A]}{\text{E}[B]^3}\text{Var}\left[B\right]
+\text{E}\left[R_2\right]
\]

From the warm-up calculations, we have: \[
\begin{aligned}
\text{E}\left[A\right] &= 0 \\
\text{E}\left[B\right] &= \frac{\sigma^2}{1-\rho^{2}}
\end{aligned}
\]

Substituting these into the expansion simplifies the expression
considerably.

First, the leading term vanishes: \[
\frac{\text{E}[A]}{\text{E}[B]} = 0
\]

Similarly, the variance correction term disappears: \[
\frac{\text{E}[A]}{\text{E}[B]^3}\text{Var}\left[B\right] = 0
\]

This leaves only the covariance term: \[
\text{E}\left[\frac{A}{B}\right]
=-\frac{1}{\text{E}[B]^2}\text{Cov}\left[A,B\right]
+\text{E}\left[R_2\right]
\]

\begin{center}\rule{0.5\linewidth}{0.5pt}\end{center}

\subsubsection{Substituting known
quantities}\label{substituting-known-quantities}

We now plug in: - \(\text{E}[B] = \frac{\sigma^2}{1-\rho^2}\) - the
previously derived expression for \(\text{Cov}(A,B)\)

Because \(\text{E}[B]^2 = \frac{\sigma^4}{(1-\rho^2)^2}\), the scaling
simplifies neatly when combined with \(\text{Cov}(A,B)\).

Carrying out this substitution and simplifying yields: \[
\text{E}\left[\frac{A}{B}\right]
= -\frac{2\rho}{T-1} + O(T^{-2})
\]

\begin{center}\rule{0.5\linewidth}{0.5pt}\end{center}

\subsubsection{Final result}\label{final-result}

Recalling that \[
\hat\beta_1 - \rho = \frac{A}{B},
\] we obtain the bias of the estimator: \[
\text{E}\left[\hat\beta_{1}\right]-\rho
= -\frac{2\rho}{T-1}+O(T^{-2})
\]

\begin{center}\rule{0.5\linewidth}{0.5pt}\end{center}

\subsubsection{Interpretation}\label{interpretation}

This is the classical Kendall bias result for the AR(1) model without a
constant.

Key takeaways: - The bias is of order \(O(1/T)\), so it vanishes
asymptotically - The bias is \textbf{negative when \(\rho > 0\)},
meaning the estimator systematically underestimates persistence in
finite samples - The magnitude of the bias increases with \(\rho\), so
it is most severe for highly persistent processes

Conceptually, the bias arises entirely from the covariance between the
numerator and denominator in the OLS estimator. Even though the
numerator has mean zero, its interaction with the random denominator
produces a systematic distortion.

This completes the derivation: we have shown how the ratio structure of
the estimator, combined with dependence in the data, generates
finite-sample bias.

    \section{Full problem}\label{full-problem}

We now consider the more general AR(1) model with an intercept: \[
y_t = \alpha+\rho y_{t-1}+e_t
\]

In regression form, this becomes \[
y_t = \beta_0+\beta_1 y_{t-1}+u_t
\]

where under the true model, \(\beta_0 = \alpha\), \(\beta_1 = \rho\),
and \(u_t = e_t\).

\begin{center}\rule{0.5\linewidth}{0.5pt}\end{center}

\subsubsection{Why this case is
different}\label{why-this-case-is-different}

Unlike the no-constant case, the process now has a nonzero mean. In
particular, under stationarity, \[
\text{E}[y_t] = \frac{\alpha}{1-\rho}
\]

Because of this, OLS must estimate both an intercept and a slope. This
changes the form of the estimator: instead of simple sums, everything is
expressed in terms of \textbf{demeaned variables}.

\begin{center}\rule{0.5\linewidth}{0.5pt}\end{center}

\subsubsection{OLS estimator with a
constant}\label{ols-estimator-with-a-constant}

The OLS estimator for \(\beta_1\) is \[
\begin{aligned}
\hat\beta_{1}
&= \frac{\sum_{t=2}^{T} (y_t-\bar y_0)(y_{t-1}-\bar y_{-1})}
{\sum_{t=2}^{T} (y_{t-1}-\bar y_{-1})^2}
\end{aligned}
\]

where: - \(\bar y_0 = (T-1)^{-1}\sum_{t=2}^T y_t\) -
\(\bar y_{-1} = (T-1)^{-1}\sum_{t=2}^T y_{t-1}\)

These are the sample means of the dependent variable and the lagged
regressor.

\begin{center}\rule{0.5\linewidth}{0.5pt}\end{center}

\subsubsection{Substituting the true
model}\label{substituting-the-true-model}

We now substitute the data generating process into the numerator: \[
\begin{aligned}
\hat\beta_{1}
&= \frac{\sum_{t=2}^{T} (\alpha+\rho y_{t-1}+e_t-\bar y_0)(y_{t-1}-\bar y_{-1})}
{\sum_{t=2}^{T} (y_{t-1}-\bar y_{-1})^2}
\end{aligned}
\]

Next, we expand \(\bar y_0\). Since \[
\bar y_0 = \alpha + \rho \bar y_{-1} + \bar e_0,
\] we can rewrite the numerator as \[
\begin{aligned}
(\alpha+\rho y_{t-1}+e_t-\alpha-\rho \bar y_{-1}-\bar e_0)
\end{aligned}
\]

This gives \[
\begin{aligned}
\hat\beta_{1}
&= \frac{\sum_{t=2}^{T} (\alpha+\rho y_{t-1}+e_t-\alpha-\rho \bar y_{-1}-\bar e_0)(y_{t-1}-\bar y_{-1})}
{\sum_{t=2}^{T} (y_{t-1}-\bar y_{-1})^2}
\end{aligned}
\]

\begin{center}\rule{0.5\linewidth}{0.5pt}\end{center}

\subsubsection{Simplification}\label{simplification}

The constant terms cancel, leaving only deviations from the mean: \[
\alpha - \alpha = 0
\]

We are left with \[
\begin{aligned}
\hat\beta_{1}
&= \frac{\sum_{t=2}^{T} (\rho(y_{t-1}-\bar y_{-1}) + (e_t - \bar e_0))(y_{t-1}-\bar y_{-1})}
{\sum_{t=2}^{T} (y_{t-1}-\bar y_{-1})^2}
\end{aligned}
\]

Expanding the numerator: \[
\begin{aligned}
&= \frac{\rho \sum_{t=2}^{T} (y_{t-1}-\bar y_{-1})^2
+ \sum_{t=2}^{T} (e_t - \bar e_0)(y_{t-1}-\bar y_{-1})}
{\sum_{t=2}^{T} (y_{t-1}-\bar y_{-1})^2}
\end{aligned}
\]

The first term simplifies directly, yielding \[
\hat\beta_{1}
= \rho+\frac{\sum_{t=2}^{T} (e_t - \bar e_0)(y_{t-1}-\bar y_{-1})}
{\sum_{t=2}^{T} (y_{t-1}-\bar y_{-1})^2}
\]

Finally, using the fact that \(\sum (y_{t-1}-\bar y_{-1}) = 0\), the
\(\bar e_0\) term drops out: \[
\sum_{t=2}^T \bar e_0 (y_{t-1}-\bar y_{-1}) = 0
\]

This gives the clean decomposition \[
\hat\beta_{1}
= \rho+\frac{\sum_{t=2}^{T} e_t(y_{t-1}-\bar y_{-1})}
{\sum_{t=2}^{T} (y_{t-1}-\bar y_{-1})^2}
\]

\begin{center}\rule{0.5\linewidth}{0.5pt}\end{center}

\subsubsection{Interpretation}\label{interpretation}

This expression is the direct analogue of the no-constant case, but with
an important difference:

\begin{itemize}
\tightlist
\item
  In the no-constant case, the regressor is \(y_{t-1}\)
\item
  With a constant, the regressor becomes \((y_{t-1} - \bar y_{-1})\)
\end{itemize}

This centering introduces additional dependence through the sample mean
\(\bar y_{-1}\), which itself depends on all observations.

As a result: - the numerator and denominator become more complicated
objects - additional covariance terms appear - the bias calculation
becomes significantly more involved

However, structurally, the problem is the same: we again have a ratio of
two random quantities, and we will analyze it using the same Taylor
expansion framework.

    \section{Computations}\label{computations}

We now translate the full problem with a constant into the same
analytical framework used in the no-constant case.

From the previous derivation, we have the decomposition \[
\hat\beta_{1}-\rho
= \frac{(T-1)^{-1}\sum_{t=2}^{T} e_t(y_{t-1}-\bar y_{-1})}
{(T-1)^{-1}\sum_{t=2}^{T} (y_{t-1}-\bar y_{-1})^2}
= \frac{A}{B}
\]

where we define \[
A = (T-1)^{-1}\sum_{t=2}^{T} e_t(y_{t-1}-\bar y_{-1})
\] \[
B = (T-1)^{-1}\sum_{t=2}^{T} (y_{t-1}-\bar y_{-1})^2
\]

\begin{center}\rule{0.5\linewidth}{0.5pt}\end{center}

\subsubsection{Strategy}\label{strategy}

Our objective is to compute \[
\text{E}[\hat\beta_1 - \rho] = \text{E}\left[\frac{A}{B}\right]
\]

As before, this is the expectation of a ratio of random variables. We
handle this using the same Taylor expansion approximation: \[
\text{E}\left[\frac{A}{B}\right]
=\frac{\text{E}[A]}{\text{E}[B]}
-\frac{1}{\text{E}[B]^2}\text{Cov}\left[A,B\right]
+\frac{\text{E}[A]}{\text{E}[B]^3}\text{Var}\left[B\right]
+\text{E}\left[R_2\right]
\]

\begin{center}\rule{0.5\linewidth}{0.5pt}\end{center}

\subsubsection{What changes relative to the no-constant
case}\label{what-changes-relative-to-the-no-constant-case}

At first glance, this looks identical to the earlier setup. However, the
presence of the sample mean \(\bar y_{-1}\) fundamentally changes the
structure of both \(A\) and \(B\):

\begin{itemize}
\tightlist
\item
  \(A\) now involves \(y_{t-1} - \bar y_{-1}\) instead of \(y_{t-1}\)
\item
  \(B\) is the sample variance of \(y_{t-1}\) rather than a simple
  second moment
\item
  \(\bar y_{-1}\) depends on the entire sample, introducing global
  dependence across all terms
\end{itemize}

This has several important consequences:

\begin{enumerate}
\def\labelenumi{\arabic{enumi}.}
\item
  \textbf{Loss of simple independence structure}\\
  In the no-constant case, many terms vanished because \(e_t\) was
  independent of past values.\\
  Here, \(\bar y_{-1}\) depends on all observations, including those
  involving \(e_t\), so fewer terms drop out.
\item
  \textbf{Additional covariance channels}\\
  The sample mean introduces new interactions between observations at
  different time indices, which will appear in \(\text{Cov}(A,B)\).
\item
  \textbf{More complex algebra, same principle}\\
  Despite the added complexity, the overall strategy remains unchanged:

  \begin{itemize}
  \tightlist
  \item
    compute \(\text{E}[A]\)
  \item
    compute \(\text{E}[B]\)
  \item
    compute \(\text{Cov}(A,B)\)
  \item
    plug into the Taylor expansion
  \end{itemize}
\end{enumerate}

\begin{center}\rule{0.5\linewidth}{0.5pt}\end{center}

\subsubsection{Roadmap}\label{roadmap}

The remainder of the derivation will proceed exactly as before, but with
more involved calculations:

\begin{itemize}
\tightlist
\item
  First, we compute \(\text{E}[A]\) and \(\text{E}[B]\), paying careful
  attention to the role of \(\bar y_{-1}\)
\item
  Next, we evaluate \(\text{Cov}(A,B)\), where most of the additional
  complexity arises
\item
  Finally, we combine these pieces to obtain the bias
\end{itemize}

Conceptually, nothing new is happening---the bias still comes from the
interaction between the numerator and denominator---but the presence of
the sample mean makes those interactions much richer.

    \subsection{\texorpdfstring{Compute
\(\text{E}[A]\)}{Compute \textbackslash text\{E\}{[}A{]}}}\label{compute-textea}

We now compute the expectation of \[
A = (T-1)^{-1}\sum_{t=2}^{T} e_t(y_{t-1}-\bar y_{-1})
\]

\begin{center}\rule{0.5\linewidth}{0.5pt}\end{center}

\subsubsection{Step 1: Expand the
definition}\label{step-1-expand-the-definition}

Taking expectations, \[
\begin{aligned}
\text{E}[A]
&= (T-1)^{-1}\sum_{t=2}^{T} \text{E}\left[e_t(y_{t-1}-\bar y_{-1})\right]\\
&= (T-1)^{-1}\sum_{t=2}^{T} \left(\text{E}[e_t y_{t-1}] - \text{E}[e_t \bar y_{-1}]\right)
\end{aligned}
\]

In the no-constant case, both terms would be zero. However, the presence
of \(\bar y_{-1}\) changes things.

\begin{itemize}
\item
  The first term is still zero: \[
  \text{E}[e_t y_{t-1}] = 0
  \] because \(e_t\) is independent of past values.
\item
  The second term is \textbf{not zero}, because \(\bar y_{-1}\) contains
  future observations relative to \(t\), which depend on \(e_t\).
\end{itemize}

\begin{center}\rule{0.5\linewidth}{0.5pt}\end{center}

\subsubsection{Step 2: Expand the sample
mean}\label{step-2-expand-the-sample-mean}

Recall that \[
\bar y_{-1} = (T-1)^{-1}\sum_{s=2}^{T} y_{s-1}
\]

Substituting this in, \[
\begin{aligned}
\text{E}[A]
&= -(T-1)^{-2}\sum_{t=2}^{T}\sum_{s=2}^{T} \text{E}[e_t y_{s-1}]
\end{aligned}
\]

\begin{center}\rule{0.5\linewidth}{0.5pt}\end{center}

\subsubsection{Step 3: Identify nonzero
terms}\label{step-3-identify-nonzero-terms}

We now analyze when \(\text{E}[e_t y_{s-1}]\) is nonzero.

\begin{itemize}
\tightlist
\item
  If \(s \le t\), then \(y_{s-1}\) depends only on shocks up to time
  \(s-1 \le t-1\), which are independent of \(e_t\), so the expectation
  is zero.
\item
  If \(s > t\), then \(y_{s-1}\) contains \(e_t\) in its expansion, so
  the expectation is nonzero.
\end{itemize}

Thus, we restrict the sum to \(s > t\): \[
\begin{aligned}
\text{E}[A]
&= -\frac{1}{(T-1)^2}\sum_{t=2}^{T}\sum_{s=t+1}^{T} \text{E}[e_t y_{s-1}]
\end{aligned}
\]

\begin{center}\rule{0.5\linewidth}{0.5pt}\end{center}

\subsubsection{Step 4: Use the AR(1)
representation}\label{step-4-use-the-ar1-representation}

Using the moving average form, \[
y_{s-1} = \sum_{q=0}^{s-2}\rho^q e_{s-1-q} + \rho^{s-1}y_0
\]

Only the shock terms contribute, since \(\text{E}[e_t y_0] = 0\).
Substituting, \[
\begin{aligned}
\text{E}[A]
&= -\frac{1}{(T-1)^2}\sum_{t=2}^{T}\sum_{s=t+1}^{T}
\text{E}\left[e_t \sum_{q=0}^{s-2}\rho^q e_{s-1-q}\right]
\end{aligned}
\]

\begin{center}\rule{0.5\linewidth}{0.5pt}\end{center}

\subsubsection{Step 5: Match indices}\label{step-5-match-indices}

The only nonzero contributions occur when the shocks match: \[
e_t \cdot e_{s-1-q} \neq 0 \quad \text{only if} \quad t = s-1-q
\]

This implies \(q = s-1-t\), so the corresponding term is
\(\rho^{s-1-t}\).

Using \(\text{E}[e_t^2] = \sigma^2\), we obtain \[
\begin{aligned}
\text{E}[A]
&= -\frac{\sigma^2}{(T-1)^2}\sum_{t=2}^{T}\sum_{s=t+1}^{T} \rho^{s-1-t}
\end{aligned}
\]

\begin{center}\rule{0.5\linewidth}{0.5pt}\end{center}

\subsubsection{Step 6: Evaluate the geometric
sum}\label{step-6-evaluate-the-geometric-sum}

For fixed \(t\), the inner sum is geometric: \[
\sum_{s=t+1}^{T} \rho^{s-1-t}
= \sum_{k=0}^{T-t-1} \rho^k
= \frac{1-\rho^{T-t}}{1-\rho}
\]

Substituting, \[
\begin{aligned}
\text{E}[A]
&= -\frac{\sigma^2}{(T-1)^2}\sum_{t=2}^{T} \frac{1-\rho^{T-t}}{1-\rho}
\end{aligned}
\]

\begin{center}\rule{0.5\linewidth}{0.5pt}\end{center}

\subsubsection{Step 7: Final
simplification}\label{step-7-final-simplification}

Summing over \(t\) gives \[
\begin{aligned}
\text{E}[A]
&= -\frac{\sigma^2}{(T-1)^2(1-\rho)}
\left[(T-1) - \frac{1-\rho^{T-1}}{1-\rho}\right]
\end{aligned}
\]

\begin{center}\rule{0.5\linewidth}{0.5pt}\end{center}

\subsubsection{Interpretation}\label{interpretation}

This is the first major departure from the no-constant case:

\begin{itemize}
\tightlist
\item
  Previously, \(\text{E}[A] = 0\)
\item
  Here, \(\text{E}[A] \neq 0\)
\end{itemize}

The reason is entirely due to the sample mean \(\bar y_{-1}\): - it
introduces dependence between \(e_t\) and future observations - this
creates nonzero contributions when \(s > t\)

Importantly, \(\text{E}[A]\) is of order \(O(1/T)\), which means it
contributes directly to the leading-order bias.

This is one of the key sources of the additional bias introduced by
estimating the intercept.

    \subsection{\texorpdfstring{Compute
\(\text{E}[B]\)}{Compute \textbackslash text\{E\}{[}B{]}}}\label{compute-texteb}

We now compute the expectation of \[
B = (T-1)^{-1}\sum_{t=2}^{T} (y_{t-1}-\bar y_{-1})^2
\]

This quantity is the \textbf{sample variance} of \(y_{t-1}\). Unlike the
no-constant case, this is no longer just a second moment---it includes a
correction for the sample mean.

\begin{center}\rule{0.5\linewidth}{0.5pt}\end{center}

\subsubsection{Step 1: Expand the
square}\label{step-1-expand-the-square}

Taking expectations, \[
\begin{aligned}
\text{E}[B]
&= (T-1)^{-1}\sum_{t=2}^{T} \text{E}\big[(y_{t-1}-\bar y_{-1})^2\big]\\
&= (T-1)^{-1}\sum_{t=2}^{T} \left(\text{E}[y_{t-1}^2] - 2\text{E}[y_{t-1}\bar y_{-1}] + \text{E}[\bar y_{-1}^2]\right)
\end{aligned}
\]

At this point, everything hinges on understanding how the sample mean
interacts with individual observations.

\begin{center}\rule{0.5\linewidth}{0.5pt}\end{center}

\subsubsection{Step 2: Use stationarity and
symmetry}\label{step-2-use-stationarity-and-symmetry}

Because the process is stationary: - \(\text{E}[y_{t-1}^2] = \gamma_0\)
- \(\text{E}[y_{t-1}y_{s-1}] = \gamma_{|t-s|}\)

Also, by symmetry, each \(y_{t-1}\) has the same relationship with
\(\bar y_{-1}\).

This leads to a key simplification: \[
\begin{aligned}
\text{E}[B]
&= \gamma_0 - \text{E}[\bar y_{-1}^2]
\end{aligned}
\]

This is an important result:\\
the expected sample variance equals the true variance minus the variance
of the sample mean.

\begin{center}\rule{0.5\linewidth}{0.5pt}\end{center}

\subsubsection{\texorpdfstring{Step 3: Compute
\(\text{E}[\bar y_{-1}^2]\)}{Step 3: Compute \textbackslash text\{E\}{[}\textbackslash bar y\_\{-1\}\^{}2{]}}}\label{step-3-compute-textebar-y_-12}

Recall that \[
\bar y_{-1} = (T-1)^{-1}\sum_{t=2}^{T} y_{t-1}
\]

Squaring and taking expectations, \[
\begin{aligned}
\text{E}[\bar y_{-1}^2]
&= (T-1)^{-2}\sum_{t=2}^{T}\sum_{s=2}^{T} \text{E}[y_{t-1}y_{s-1}]
\end{aligned}
\]

Using stationarity, \[
\begin{aligned}
&= (T-1)^{-2}\sum_{t=2}^{T}\sum_{s=2}^{T} \gamma_{|t-s|}
\end{aligned}
\]

We now reorganize the double sum by lag \(h = |t-s|\): \[
\begin{aligned}
&= (T-1)^{-2}\left[(T-1)\gamma_0 + 2\sum_{h=1}^{T-2}(T-1-h)\gamma_h\right]
\end{aligned}
\]

This expression reflects: - \((T-1)\) terms with lag 0 - \((T-1-h)\)
pairs at each lag \(h\), counted twice for symmetry

\begin{center}\rule{0.5\linewidth}{0.5pt}\end{center}

\subsubsection{Step 4: Substitute AR(1)
autocovariances}\label{step-4-substitute-ar1-autocovariances}

For the AR(1) process, \[
\gamma_h = \frac{\sigma^2}{1-\rho^2}\rho^h
\]

Substituting into \(\text{E}[B]\) gives \[
\text{E}[B]
=
\frac{\sigma^2}{1-\rho^2}
-
(T-1)^{-2}
\left[
(T-1)\frac{\sigma^2}{1-\rho^2}
+
2\sum_{h=1}^{T-2}(T-1-h)\frac{\sigma^2}{1-\rho^2}\rho^h
\right]
\]

Factoring out \(\frac{\sigma^2}{1-\rho^2}\): \[
\text{E}[B]
=
\frac{\sigma^2}{1-\rho^2}
\left[
1
-
\frac{1}{(T-1)^2}
\left(
(T-1)
+
2\sum_{h=1}^{T-2}(T-1-h)\rho^h
\right)
\right]
\]

\begin{center}\rule{0.5\linewidth}{0.5pt}\end{center}

\subsubsection{Step 5: Evaluate the weighted geometric
sum}\label{step-5-evaluate-the-weighted-geometric-sum}

The remaining sum is a weighted geometric series: \[
\sum_{h=1}^{T-2}(T-1-h)\rho^h
\]

Evaluating this yields \[
\begin{aligned}
&=
(T-1)\frac{\rho(1-\rho^{T-2})}{1-\rho}
-
\frac{\rho\big(1-(T-1)\rho^{T-2}+(T-2)\rho^{T-1}\big)}{(1-\rho)^2}
\end{aligned}
\]

Substituting back, \[
\begin{aligned}
\text{E}[B]
&=
\frac{\sigma^2}{1-\rho^2}
-
\frac{\sigma^2}{(T-1)^2(1-\rho^2)}
\Bigg[
(T-1)
+
2\Bigg(
(T-1)\frac{\rho(1-\rho^{T-2})}{1-\rho}\\
&\qquad\qquad\qquad
-
\frac{\rho\big(1-(T-1)\rho^{T-2}+(T-2)\rho^{T-1}\big)}{(1-\rho)^2}
\Bigg)
\Bigg]
\end{aligned}
\]

\begin{center}\rule{0.5\linewidth}{0.5pt}\end{center}

\subsubsection{Step 6: Asymptotic
simplification}\label{step-6-asymptotic-simplification}

Although the exact expression is complex, its asymptotic behavior is
simple. The correction term is of order \(O(1/T)\), so \[
\text{E}[B] = \frac{\sigma^2}{1-\rho^2} + O\!\left(T^{-1}\right)
\]

\begin{center}\rule{0.5\linewidth}{0.5pt}\end{center}

\subsubsection{Interpretation}\label{interpretation}

This result shows that:

\begin{itemize}
\tightlist
\item
  The expected denominator is close to the unconditional variance
  \(\gamma_0\)
\item
  The correction comes from estimating the mean using the same data
\item
  This correction is small (order \(1/T\)), but it plays a role in the
  finite-sample bias
\end{itemize}

Conceptually: - In the no-constant case, \(B\) directly estimates
\(\gamma_0\) - With a constant, \(B\) estimates \(\gamma_0\)
\textbf{minus the variability of the sample mean}

This is another way the intercept complicates the problem, introducing
additional small-sample effects that must be accounted for in the bias
calculation.

    \subsection{\texorpdfstring{Compute \(\text{Var}[B]\) (nuisance
term)}{Compute \textbackslash text\{Var\}{[}B{]} (nuisance term)}}\label{compute-textvarb-nuisance-term}

We now compute \(\text{Var}(B)\), which appears in the Taylor
expansion.\\
This term will ultimately be of smaller importance, but we include it
for completeness.

\begin{center}\rule{0.5\linewidth}{0.5pt}\end{center}

\subsubsection{\texorpdfstring{Step 1: Rewrite
\(B\)}{Step 1: Rewrite B}}\label{step-1-rewrite-b}

Let \[
\text{Let } n=T-1,\qquad x_t=y_{t-1}
\]

Then \[
\begin{aligned}
B
&= n^{-1}\sum_{t=1}^n (x_t-\bar x)^2\\
&= n^{-1}\sum_{t=1}^n x_t^2-\bar x^2
\end{aligned}
\]

This decomposition separates \(B\) into: - the average of squared
observations - minus the square of the sample mean

\begin{center}\rule{0.5\linewidth}{0.5pt}\end{center}

\subsubsection{Step 2: Identify leading-order
behavior}\label{step-2-identify-leading-order-behavior}

We now analyze the variance: \[
\text{Var}(B)
\]

The key observation is that \(\bar x = O_p(n^{-1/2})\), so \[
\bar x^2 = O_p(n^{-1})
\]

This implies \[
\text{Var}(\bar x^2) = O(n^{-2})
\]

which is negligible relative to the \(O(n^{-1})\) scale of the bias.
Therefore, to leading order, \[
\text{Var}(B)=\text{Var}\!\left(n^{-1}\sum_{t=1}^n x_t^2\right)+o\!\left(\frac{1}{n}\right)
\]

So we can focus on the variance of the average of \(x_t^2\).

\begin{center}\rule{0.5\linewidth}{0.5pt}\end{center}

\subsubsection{Step 3: Expand the
variance}\label{step-3-expand-the-variance}

Using the standard formula for the variance of a sum, \[
\begin{aligned}
\text{Var}\!\left(n^{-1}\sum_{t=1}^n x_t^2\right)
&= n^{-2}\sum_{t=1}^n\sum_{s=1}^n \text{Cov}(x_t^2,x_s^2)
\end{aligned}
\]

By stationarity, this depends only on the lag \(h = |t-s|\), giving \[
\begin{aligned}
&= n^{-2}\left[n\,\text{Var}(x_1^2)+2\sum_{h=1}^{n-1}(n-h)\text{Cov}(x_1^2,x_{1+h}^2)\right]
\end{aligned}
\]

\begin{center}\rule{0.5\linewidth}{0.5pt}\end{center}

\subsubsection{Step 4: Use Gaussian
structure}\label{step-4-use-gaussian-structure}

For a Gaussian stationary AR(1), \[
\gamma_h=\text{Cov}(x_t,x_{t+h})=\frac{\sigma^2}{1-\rho^2}\rho^h
\]

To compute \(\text{Cov}(x_t^2,x_{t+h}^2)\), we use Isserlis' formula
(moment factoring for Gaussian variables): \[
\text{Cov}(x_t^2,x_{t+h}^2)=2\gamma_h^2
\]

This is a major simplification: fourth moments reduce to functions of
second moments.

\begin{center}\rule{0.5\linewidth}{0.5pt}\end{center}

\subsubsection{Step 5: Substitute and
simplify}\label{step-5-substitute-and-simplify}

Substituting into the variance expression, \[
\begin{aligned}
\text{Var}(B)
&= 2n^{-2}\left[n\gamma_0^2+2\sum_{h=1}^{n-1}(n-h)\gamma_h^2\right]+o\!\left(\frac{1}{n}\right)
\end{aligned}
\]

Using \(\gamma_h = \gamma_0 \rho^h\), we get \[
\begin{aligned}
\text{Var}(B)
&= 2\gamma_0^2 n^{-2}\left[n+2\sum_{h=1}^{n-1}(n-h)\rho^{2h}\right]+o\!\left(\frac{1}{n}\right)
\end{aligned}
\]

\begin{center}\rule{0.5\linewidth}{0.5pt}\end{center}

\subsubsection{Step 6: Evaluate the weighted geometric
sum}\label{step-6-evaluate-the-weighted-geometric-sum}

The remaining sum is \[
\sum_{h=1}^{n-1}(n-h)\rho^{2h}
\]

which evaluates to \[
\begin{aligned}
&=
n\frac{\rho^2(1-\rho^{2(n-1)})}{1-\rho^2}
-
\frac{\rho^2\bigl(1-n\rho^{2n-2}+(n-1)\rho^{2n}\bigr)}{(1-\rho^2)^2}
\end{aligned}
\]

Substituting this back yields the full expression: \[
\begin{aligned}
\text{Var}(B)
&=
\frac{2\gamma_0^2}{n^2}
\left[
n
+
2n\frac{\rho^2(1-\rho^{2(n-1)})}{1-\rho^2}
-
2\frac{\rho^2\bigl(1-n\rho^{2n-2}+(n-1)\rho^{2n}\bigr)}{(1-\rho^2)^2}
\right]
+o\!\left(\frac{1}{n}\right)
\end{aligned}
\]

\begin{center}\rule{0.5\linewidth}{0.5pt}\end{center}

\subsubsection{Step 7: Leading-order
behavior}\label{step-7-leading-order-behavior}

Using \(\gamma_0 = \frac{\sigma^2}{1-\rho^2}\) and simplifying, we
obtain \[
\text{Var}(B)
=
\frac{2\sigma^4(1+\rho^2)}{n(1-\rho^2)^3}
+O\!\left(n^{-2}\right)
\]

Returning to \(T\), \[
\text{Var}(B)
=
\frac{2\sigma^4(1+\rho^2)}{(T-1)(1-\rho^2)^3}
+O\!\left(T^{-2}\right)
\]

\begin{center}\rule{0.5\linewidth}{0.5pt}\end{center}

\subsubsection{Interpretation}\label{interpretation}

This term is labeled a \emph{nuisance term} because:

\begin{itemize}
\tightlist
\item
  It is multiplied by \(\text{E}[A]\) in the Taylor expansion
\item
  Since \(\text{E}[A] = O(1/T)\) and \(\text{Var}(B) = O(1/T)\), their
  product is \(O(1/T^2)\)
\item
  Therefore, it does \textbf{not} contribute to the leading-order bias
\end{itemize}

Conceptually: - \(\text{Var}(B)\) measures how much the denominator
fluctuates - But these fluctuations only matter at second order

So while this calculation is technically involved, its role in the final
result is limited: it confirms that denominator variability does not
affect the leading \(O(1/T)\) bias.

    \subsection{\texorpdfstring{Compute
\(\text{Cov}[A,B]\)}{Compute \textbackslash text\{Cov\}{[}A,B{]}}}\label{compute-textcovab}

We now compute the covariance term \[
\text{Cov}[A,B]
\] for the full problem with a constant.

Recall the definitions \[
\begin{aligned}
A &= (T-1)^{-1}\sum_{t=2}^{T} e_t(y_{t-1}-\bar y_{-1}),\\
B &= (T-1)^{-1}\sum_{t=2}^{T} (y_{t-1}-\bar y_{-1})^2.
\end{aligned}
\]

This is the most complicated part of the derivation. Conceptually,
however, it is doing the same job as before: it measures how the random
numerator and random denominator move together.

\begin{center}\rule{0.5\linewidth}{0.5pt}\end{center}

\subsubsection{Step 1: Start from the
definition}\label{step-1-start-from-the-definition}

By definition, \[
\begin{aligned}
\text{Cov}[A,B]
&=\text{E}\big[(A-\text{E}[A])(B-\text{E}[B])\big]
\end{aligned}
\]

Using the fact that, to leading order, \[
\text{E}[B] = \frac{\sigma^2}{1-\rho^2} + O(T^{-1}),
\] we write \[
\begin{aligned}
\text{Cov}[A,B]
&=\text{E}\left[A\left(B-\frac{\sigma^2}{1-\rho^2}\right)\right]
\end{aligned}
\]

Expanding the product gives \[
\begin{aligned}
\text{Cov}[A,B]
&=\text{E}\left[(T-1)^{-2}\sum_{t=2}^{T}\sum_{s=2}^{T} e_t y_{t-1}y_{s-1}^2\right].
\end{aligned}
\]

At this stage, the expression is formally the same as in the no-constant
case. The difference is that in the full problem this covariance is no
longer the only contributor to the bias, because now
\(\text{E}[A] \neq 0\) as well.

\begin{center}\rule{0.5\linewidth}{0.5pt}\end{center}

\subsubsection{Step 2: Restrict to nonzero
terms}\label{step-2-restrict-to-nonzero-terms}

As before, only terms with \(s>t\) contribute.

The reason is that if \(s \le t\), then \(y_{s-1}\) depends only on
shocks dated no later than \(t-1\), so it is independent of \(e_t\).
Hence those terms have zero expectation.

Therefore, \[
\begin{aligned}
\text{Cov}[A,B]
&=\text{E}\left[(T-1)^{-2}\sum_{t=2}^{T}\sum_{s=t+1}^{T} e_t y_{t-1}y_{s-1}^2\right].
\end{aligned}
\]

This is the same forward-dependence argument used in the warm-up
calculation.

\begin{center}\rule{0.5\linewidth}{0.5pt}\end{center}

\subsubsection{Step 3: Substitute the AR(1)
representation}\label{step-3-substitute-the-ar1-representation}

We now write both \(y_{t-1}\) and \(y_{s-1}\) in moving-average form: \[
y_{t-1}=\rho^{t-1}y_0+\sum_{r=0}^{t-2}\rho^r e_{t-1-r},
\qquad
y_{s-1}=\rho^{s-1}y_0+\sum_{q=0}^{s-2}\rho^q e_{s-1-q}.
\]

Substituting these expansions into \[
e_t y_{t-1} y_{s-1}^2
\] produces a large number of terms. Most of them vanish after taking
expectations, because they involve odd products of mean-zero shocks.

After expanding and discarding all odd-moment terms, only two types of
contributions remain:

\[
\begin{aligned}
\text{Cov}[A,B]
&=
(T-1)^{-2}\sum_{t=2}^{T}\sum_{s=t+1}^{T}\sum_{p=0}^{s-2}
2\rho^{t-1}\rho^{s-1}\rho^p\,\text{E}[y_0^2 e_t e_{s-1-p}]\\
&\quad+
(T-1)^{-2}\sum_{t=2}^{T}\sum_{s=t+1}^{T}\sum_{r=0}^{t-2}\sum_{q=0}^{s-2}\sum_{p=0}^{s-2}
\rho^r\rho^{q+p}\,\text{E}[e_t e_{t-1-r}e_{s-1-q}e_{s-1-p}].
\end{aligned}
\]

So the surviving terms are: - one term involving \(y_0^2\) and two
shocks - one term involving four shocks

\begin{center}\rule{0.5\linewidth}{0.5pt}\end{center}

\subsubsection{Step 4: Evaluate the first surviving
term}\label{step-4-evaluate-the-first-surviving-term}

Consider first \[
\begin{aligned}
&(T-1)^{-2}\sum_{t=2}^{T}\sum_{s=t+1}^{T}\sum_{p=0}^{s-2}
2\rho^{t-1}\rho^{s-1}\rho^p\,\text{E}[y_0^2 e_t e_{s-1-p}]
\end{aligned}
\]

The expectation is nonzero only when the shocks match: \[
s-1-p=t
\qquad\Longrightarrow\qquad
p=s-1-t.
\]

Substituting this matching index, we obtain \[
\begin{aligned}
&(T-1)^{-2}\sum_{t=2}^{T}\sum_{s=t+1}^{T}
2\rho^{t-1}\rho^{s-1}\rho^{s-1-t}\,\text{E}[y_0^2]\sigma^2
\end{aligned}
\]

Using stationarity, \[
\text{E}[y_0^2]=\frac{\sigma^2}{1-\rho^2},
\] so this becomes \[
\begin{aligned}
&=
(T-1)^{-2}\sum_{t=2}^{T}\sum_{s=t+1}^{T}
2\rho^{2s-3}(1-\rho^2)^{-1}\sigma^4.
\end{aligned}
\]

\begin{center}\rule{0.5\linewidth}{0.5pt}\end{center}

\subsubsection{Step 5: Evaluate the four-shock
term}\label{step-5-evaluate-the-four-shock-term}

Now consider \[
\begin{aligned}
&(T-1)^{-2}\sum_{t=2}^{T}\sum_{s=t+1}^{T}\sum_{r=0}^{t-2}\sum_{q=0}^{s-2}\sum_{p=0}^{s-2}
\rho^r\rho^{q+p}\,\text{E}[e_t e_{t-1-r}e_{s-1-q}e_{s-1-p}].
\end{aligned}
\]

The expectation is nonzero only when the four shocks can be paired into
matching indices. There are two such nonzero pairings, which is why a
factor of 2 appears.

After matching indices, this reduces to \[
\begin{aligned}
&=
(T-1)^{-2}\sum_{t=2}^{T}\sum_{s=t+1}^{T}\sum_{r=0}^{t-2}
2\rho^r\rho^{s-1-t}\rho^{s-t+r}\sigma^4
\end{aligned}
\]

Combining exponents yields \[
\begin{aligned}
&=
(T-1)^{-2}\sum_{t=2}^{T}\sum_{s=t+1}^{T}\sum_{r=0}^{t-2}
2\rho^{2r+2s-2t-1}\sigma^4.
\end{aligned}
\]

\begin{center}\rule{0.5\linewidth}{0.5pt}\end{center}

\subsubsection{Step 6: Combine the two
pieces}\label{step-6-combine-the-two-pieces}

Putting the two surviving terms together gives \[
\begin{aligned}
\text{Cov}[A,B]
&=
(T-1)^{-2}\sum_{t=2}^{T}\sum_{s=t+1}^{T}
2\rho^{2s-3}(1-\rho^2)^{-1}\sigma^4\\
&\quad+
(T-1)^{-2}\sum_{t=2}^{T}\sum_{s=t+1}^{T}\sum_{r=0}^{t-2}
2\rho^{2r+2s-2t-1}\sigma^4.
\end{aligned}
\]

The first term is already in usable form. The second still contains an
inner geometric sum.

\begin{center}\rule{0.5\linewidth}{0.5pt}\end{center}

\subsubsection{Step 7: Evaluate the inner geometric
sum}\label{step-7-evaluate-the-inner-geometric-sum}

For the second term, \[
\begin{aligned}
\sum_{r=0}^{t-2}\rho^{2r+2s-2t-1}
&=
\rho^{2s-2t-1}\sum_{r=0}^{t-2}\rho^{2r}\\
&=
\rho^{2s-2t-1}\frac{1-\rho^{2(t-1)}}{1-\rho^2}.
\end{aligned}
\]

Substituting this back yields \[
\begin{aligned}
\text{Cov}[A,B]
&=
2(T-1)^{-2}(1-\rho^2)^{-1}\sigma^4
\left[
\sum_{t=2}^{T}\sum_{s=t+1}^{T}\rho^{2s-3}
+
\sum_{t=2}^{T}\sum_{s=t+1}^{T}\rho^{2s-2t-1}\left(1-\rho^{2(t-1)}\right)
\right].
\end{aligned}
\]

At this point, everything has been reduced to geometric series.

\begin{center}\rule{0.5\linewidth}{0.5pt}\end{center}

\subsubsection{Step 8: Closed form and
asymptotics}\label{step-8-closed-form-and-asymptotics}

Evaluating the remaining sums gives the exact closed form \[
\begin{aligned}
\text{Cov}[A,B]
&=
\frac{2\sigma^4\left(T\rho^4-T\rho^2-\rho^4+2\rho^2-\rho^{2T}\right)}
{\rho (T-1)^2(\rho-1)^3(\rho+1)^3}\\
&=
-\frac{2\sigma^4\left(T\rho^4-T\rho^2-\rho^4+2\rho^2-\rho^{2T}\right)}
{\rho (T-1)^2(1-\rho^2)^3}.
\end{aligned}
\]

Taking the leading term as \(T\to\infty\), \[
\begin{aligned}
\text{Cov}[A,B]
&=
\frac{2\rho\,\sigma^4}{(T-1)(1-\rho^2)^2}
+O\left(T^{-2}\right).
\end{aligned}
\]

\begin{center}\rule{0.5\linewidth}{0.5pt}\end{center}

\subsubsection{Interpretation}\label{interpretation}

This shows that \(\text{Cov}[A,B]\) remains of order \(O(1/T)\) in the
full problem.

So, just as in the no-constant case: - the numerator and denominator are
not independent - their covariance contributes directly to the
finite-sample bias - the leading contribution is again at order \(1/T\)

The difference now is that this covariance term must be combined with
the nonzero value of \(\text{E}[A]\). In the no-constant case,
\(\text{E}[A]=0\), so covariance alone drove the bias. Here, both
effects matter.

    \section{Assemble the bias}\label{assemble-the-bias}

\[
\begin{aligned}
\text{E}\left[\frac{A}{B}\right]
&=\frac{\text{E}[A]}{\text{E}[B]}
-\frac{1}{\text{E}[B]^2}\text{Cov}\left[A,B\right]
+\frac{\text{E}[A]}{\text{E}[B]^3}\text{Var}\left[B\right]
+\text{E}\left[R_2\right]\\[6pt]
&=
\frac{
-\frac{\sigma^2}{(T-1)^2(1-\rho)}
\left[(T-1) - \frac{1-\rho^{T-1}}{1-\rho}\right]
}{
\frac{\sigma^2}{1-\rho^2}
}\\
&\quad
-\frac{1}{\left[\frac{\sigma^2}{1-\rho^2}\right]^2}
\cdot
\frac{2\rho\,\sigma^4}{(T-1)(1-\rho^2)^2}\\
&\quad
+\frac{
-\frac{\sigma^2}{(T-1)^2(1-\rho)}
\left[(T-1) - \frac{1-\rho^{T-1}}{1-\rho}\right]
}{
\left[\frac{\sigma^2}{1-\rho^2}\right]^3
}
\cdot
\frac{2\sigma^4(1+\rho^2)}{(T-1)(1-\rho^2)^3}
+O\!\left(T^{-2}\right)\\[6pt]
&=-\frac{1+\rho}{T-1}-\frac{2\rho}{T-1}
+O\!\left(T^{-2}\right)+O\!\left(T^{-2}\right)\\
&=-\frac{1+3\rho}{T-1}+O\!\left(T^{-2}\right)
\end{aligned}
\]

    \section{Key insights}\label{key-insights}

We now summarize the main results and explain where the bias comes from.

\begin{center}\rule{0.5\linewidth}{0.5pt}\end{center}

\subsubsection{1. Final bias formulas}\label{final-bias-formulas}

For the AR(1) model without a constant: \[
\text{E}[\hat\beta_1] - \rho
= -\frac{2\rho}{T-1} + O(T^{-2})
\]

For the AR(1) model with a constant: \[
\text{E}[\hat\beta_1] - \rho
= -\frac{1+3\rho}{T-1} + O(T^{-2})
\]

Both biases are of order \(O(1/T)\) and vanish asymptotically, but they
differ in magnitude due to the presence of the intercept.

\begin{center}\rule{0.5\linewidth}{0.5pt}\end{center}

\subsubsection{2. Where the bias comes
from}\label{where-the-bias-comes-from}

In both cases, the estimator can be written as \[
\hat\beta_1 - \rho = \frac{A}{B}
\]

The key object is therefore \[
\text{E}\!\left[\frac{A}{B}\right]
\]

Using the Taylor expansion, \[
\text{E}\!\left[\frac{A}{B}\right]
=
\frac{\text{E}[A]}{\text{E}[B]}
-\frac{1}{\text{E}[B]^2}\text{Cov}(A,B)
+\frac{\text{E}[A]}{\text{E}[B]^3}\text{Var}(B)
+O(T^{-2})
\]

This decomposition shows that bias can come from three sources: - a
nonzero numerator mean \(\text{E}[A]\) - dependence between numerator
and denominator \(\text{Cov}(A,B)\) - variability of the denominator
\(\text{Var}(B)\)

\begin{center}\rule{0.5\linewidth}{0.5pt}\end{center}

\subsubsection{3. No-constant case: bias from covariance
only}\label{no-constant-case-bias-from-covariance-only}

In the no-constant case, \[
\text{E}[A] = 0
\]

so the bias reduces to \[
\text{E}[\hat\beta_1 - \rho]
=
-\frac{1}{\text{E}[B]^2}\text{Cov}(A,B)
+O(T^{-2})
\]

Thus, the bias is driven entirely by the interaction between numerator
and denominator.

\textbf{Interpretation:}\\
Even though \(e_t\) is uncorrelated with \(y_{t-1}\) at each point in
time, the ratio structure introduces dependence between the sample
covariance (numerator) and sample variance (denominator).

\begin{center}\rule{0.5\linewidth}{0.5pt}\end{center}

\subsubsection{4. With constant: two sources of
bias}\label{with-constant-two-sources-of-bias}

With a constant, a new effect appears: \[
\text{E}[A] \neq 0
\]

so the leading terms become \[
\text{E}[\hat\beta_1 - \rho]
=
\frac{\text{E}[A]}{\text{E}[B]}
-\frac{1}{\text{E}[B]^2}\text{Cov}(A,B)
+O(T^{-2})
\]

Now the bias has two components: 1. \textbf{Mean effect} from
\(\text{E}[A]\) 2. \textbf{Covariance effect} from \(\text{Cov}(A,B)\)

This is why the coefficient changes from \(2\rho\) to \(1+3\rho\).

\begin{center}\rule{0.5\linewidth}{0.5pt}\end{center}

\subsubsection{5. Role of the sample
mean}\label{role-of-the-sample-mean}

The fundamental difference between the two cases is the presence of the
sample mean: \[
y_{t-1} \quad \longrightarrow \quad (y_{t-1} - \bar y_{-1})
\]

This introduces global dependence: - \(\bar y_{-1}\) depends on all
observations - each \(e_t\) affects future values of \(\bar y_{-1}\)

This creates terms like \[
\text{E}[e_t y_{s-1}] \neq 0 \quad \text{for } s>t
\]

which is exactly what produced \(\text{E}[A] \neq 0\).

\begin{center}\rule{0.5\linewidth}{0.5pt}\end{center}

\subsubsection{6. Order of magnitude}\label{order-of-magnitude}

All key quantities scale as: \[
\text{E}[A] = O(T^{-1}), \quad
\text{Cov}(A,B) = O(T^{-1}), \quad
\text{Var}(B) = O(T^{-1})
\]

Thus, \[
\text{E}[\hat\beta_1 - \rho] = O(T^{-1})
\]

and higher-order terms are \(O(T^{-2})\).

\begin{center}\rule{0.5\linewidth}{0.5pt}\end{center}

\subsubsection{7. Direction of the bias}\label{direction-of-the-bias}

For \(\rho > 0\): - both \(-2\rho\) and \(-(1+3\rho)\) are negative

So in both cases: \[
\text{E}[\hat\beta_1] < \rho
\]

\textbf{Interpretation:}\\
OLS underestimates persistence in finite samples.

This effect becomes more severe as \(\rho \to 1\), which is why
near-unit-root processes are especially prone to bias.

\begin{center}\rule{0.5\linewidth}{0.5pt}\end{center}

\subsubsection{8. Big picture}\label{big-picture}

The key insight is that:

\begin{quote}
\textbf{The bias is not caused by correlation between \(e_t\) and
\(y_{t-1}\), but by the randomness of the denominator.}
\end{quote}

More precisely: - OLS is estimating a ratio of random quantities - even
if the numerator has mean zero, the ratio does not - dependence across
time induces nonzero covariance between numerator and denominator

Adding a constant amplifies this effect by introducing global dependence
through the sample mean.

\begin{center}\rule{0.5\linewidth}{0.5pt}\end{center}

\subsubsection{9. Why this matters}\label{why-this-matters}

This bias has practical consequences:

\begin{itemize}
\tightlist
\item
  It explains why AR(1) estimates are systematically too small in short
  samples
\item
  It motivates bias corrections and alternative estimators
\item
  It highlights the danger of interpreting OLS estimates of persistence
  naively
\end{itemize}

In short, this is a clean example of how \textbf{finite-sample effects
arise from nonlinear transformations of random variables}, even in
otherwise simple models.


    % Add a bibliography block to the postdoc
    
    
    
\end{document}
